% ============================================================
%  sections/product.tex
% ============================================================

\section{Product Roadmap}

\subsection{MVP — Version 0.1 (Weeks 1--10)}

\textit{Goal: prove one prediction on real user data.}

\begin{itemize}
  \item Dagster sensor collecting per-run metrics into local SQLite
  \item Rolling linear trend on job latency (last 30 runs)
  \item SLA breach date calculation with confidence interval
  \item One rule: concurrency threshold → Kubernetes migration recommendation
  \item Slack bot: weekly digest + immediate alert at $\leq$14-day breach window
  \item Configuration via single \texttt{runway.yaml} (Slack token, SLA threshold)
\end{itemize}

\textbf{Not in MVP:} web dashboard, billing, multi-framework support, LLM integration.

\subsection{Version 1.0 (Months 3--6)}

\textit{Goal: five paying teams, measurable prediction accuracy.}

\begin{itemize}
  \item Airflow sensor support (expands TAM $\times$3)
  \item Prediction accuracy log (retrospective view: were predictions correct?)
  \item Web dashboard (read-only, no configuration UI yet)
  \item dbt model latency tracking (expands to analytics engineering teams)
  \item Stripe billing integration (Free / Pro / Team tiers)
\end{itemize}

\subsection{Version 2.0 (Months 7--12)}

\textit{Goal: \texteuro{}50k ARR, enterprise pipeline.}

\begin{itemize}
  \item ML pipeline support (Kubeflow, Prefect, Metaflow)
  \item LLM-generated root cause explanations (local Phi-3 via llama.cpp — privacy-preserving)
  \item Multi-pipeline correlation: "Pipeline A's failure is upstream of Pipeline B's SLA breach"
  \item SSO, RBAC, audit log — required for enterprise procurement
  \item Cloud deployment option (AWS Marketplace listing)
\end{itemize}

\subsection{Expansion path (Year 2+)}

\begin{callout}
\textbf{The natural expansion} is from pipeline orchestration to the full DataOps stack:
data warehouses (Snowflake, BigQuery query latency trends), feature stores (Feast, Tecton),
and ML training pipelines. Each expansion uses the same \emph{detect → predict → recommend}
architecture. The crossover methodology generalises: the question is always
``at what point does the current infrastructure fail this workload, and what changes that?''
\end{callout}
